\section*{Context and Motivations}
\addcontentsline{toc}{section}{Context and Motivations}

Throughout the course of its evolution, human society has been continuously and profoundly shaped by the advancement of technology.
%
From the invention of the wheel to the advent of personal computers, each technological breakthrough has had profound effects on how individuals interact with one another and with their environment.
%
Since the turn of the millennium, the rise of digital technologies has transformed the way we live, work, and communicate, propelling humankind into a new digital era.
%
As society transitioned toward this digital paradigm, the concept of data began to gain value, ultimately becoming pivotal in what is now referred to as a digital society.
%
As this digital transformation continues to unfold, driven by the widespread adoption of cloud infrastructures, smart services, and the Internet of Things (IoT), generated data has exponentially grown \mpa{{some\_ref}} not only in sheer volume but also in variety, presenting unprecedented opportunities alongside critical challenges in data management.
%
Furthermore, the availability of vast amounts of data has paved the way for the practical deployment of advanced, data-intensive technologies, such as machine learning, artificial intelligence, and digital twins.
%
Despite their long-standing theoretical foundations, these approaches remained constrained for decades by a lack of adequate computational infrastructure and sufficient data, but are now actively transforming numerous aspects of society.

Within the past decade, the emergence of digital twins has garnered significant attention across multiple domains, including manufacturing, healthcare, smart cities, and agriculture.
%
A digital twin represents one of the most advanced paradigms enabled by digital transformation, serving as a dynamic virtual replica of a physical entity, a complex system, or operational process.
%
By capturing the physical entity's underlying characteristics and continuously reflecting its state and behavior, a digital twin enables advanced simulation, predictive modeling, monitoring, and performance optimization within a virtual environment.

As reflected by their progression along the Gartner Hype Cycle~\mpa{{some\_ref}}, digital twins have evolved from an emerging research concept into an increasingly adopted technological paradigm across a wide range of application domains.
%
However, this growing maturity has also exposed significant limitations that hinder their large-scale adoption and, in particular, the development of interoperable Digital Twin.
%
Most existing digital twins have been developed as domain-specific, vertically integrated solutions, independently designed to address the requirements of individual application scenarios.
%
As a result, they rely on heterogeneous software architectures, data representations, and application-specific information models, leading to fragmented data environments in which the information generated by individual digital twins remains largely confined within isolated silos.
%
This fragmentation hinders the integration and reuse of data across independently developed digital twins, ultimately limiting interoperability, scalability, and the applicability of digital twin ecosystems to more complex, cross-domain scenarios.

Achieving interoperability is therefore becoming a fundamental requirement for the next generation of digital twin systems and, consequently, for the software infrastructures supporting them.
%
Beyond enabling information exchange, interoperability provides the foundation for integrating independently developed digital twins into larger ecosystems, where heterogeneous information can be jointly analyzed and exploited to support increasingly complex cyber-physical scenarios.
%
This challenge can be considered from two complementary dimensions.
%
The first concerns the software infrastructure supporting digital twins. Current implementations are typically developed as self-contained solutions, creating the need for common architectural foundations capable of supporting the deployment, integration, and management of heterogeneous digital twins within shared environments.
%
The second concerns the data management and semantic dimension. Even when digital twins can be deployed within a common infrastructure, the heterogeneity of the data they produce and consume, together with the use of domain-specific data models, schemas, and representations, poses significant challenges to data integration.
%
Data originating from different digital twins may differ not only in their structure and format, but also in their temporal characteristics, granularity, metadata, and underlying semantics, making their integration and joint querying a non-trivial data management problem.
%
Addressing these challenges requires common approaches to data modeling, integration, and governance, together with architectural foundations capable of supporting the heterogeneous and continuously evolving workloads associated with digital twin applications.
%
Taken together, these requirements motivate the development of software infrastructures specifically designed to support digital twin ecosystems rather than individual digital twins in isolation.

These requirements point toward the need for dedicated software infrastructures capable of supporting the deployment, management, and integration of heterogeneous digital twins within a shared environment.
%
Such infrastructures can be envisioned as Digital Twin Platforms (DTPs): software environments designed around the specific data and operational requirements of digital twin systems, providing the architectural and data management foundations required to deploy, integrate, and manage multiple digital twins within a common environment.
%
A Digital Twin Platform therefore provides a natural architectural foundation for addressing the fragmentation of current digital twin solutions, bringing together the infrastructure and data management mechanisms required to support digital twin ecosystems.
%
The importance of such foundations is likely to become even more pronounced with the recent emergence of large language models and increasingly capable agentic systems.
%
By combining access to heterogeneous information with advanced reasoning and decision-making capabilities, these technologies may enable digital twins to evolve from primarily reactive representations of physical systems toward more proactive entities capable of autonomously analyzing information, initiating actions, and, ultimately, interacting with other digital twins.
%
Emerging visions of interconnected digital twins forming distributed ``Digital Brains'' further illustrate the potential of this evolution.
%
However, realizing such scenarios requires digital twins to first overcome their current fragmentation: the ability of intelligent agents to reason over information, communicate, and coordinate actions ultimately depends on interoperable underlying systems and shared representations of the data they exchange.
%
In this context, Digital Twin Platforms can provide the common technological foundation upon which such interconnected ecosystems can be built, combining architectural support for heterogeneous digital twins with the data management capabilities required to make their information accessible, interpretable, and interoperable.

\section*{Overview and Contribution}
\addcontentsline{toc}{section}{Overview and Contribution}

This thesis approaches these challenges from a data-oriented perspective, arguing that the development of interoperable Digital Twin ecosystems should be grounded in the data they collect, produce, and exchange.
%
Rather than considering the digital twin primarily as a standalone software application, the work first investigates the characteristics and requirements of the underlying data produced and exploited by digital twins, using them as the basis for identifying suitable data models and mechanisms for their integration and management.
%
From this perspective, establishing common data models and, where possible, standardized representations for digital twins constitutes a fundamental step toward interoperability and, consequently, toward the design of effective Digital Twin Platforms.
%
Once these data requirements and representations have been established, the corresponding software infrastructure can be designed according to the characteristics of the data and the operational requirements of digital twin systems.
%
This includes supporting heterogeneous data sources, different data models and formats, varying collection frequencies, and continuous data streams, while providing the scalability and architectural capabilities required by distributed digital twin deployments.
%
The thesis therefore follows a bottom-up perspective, moving from data characteristics and requirements, through data modeling and management, toward the architectural foundations required to support interoperable Digital Twin ecosystems.
%
Within this perspective, the Digital Twin Platform represents the architectural context in which these elements are brought together: a common software foundation in which digital twins can be deployed and managed, while their heterogeneous data can be integrated, governed, and exploited across the ecosystem.
%
Such infrastructures can be envisioned as a Digital Twin Platform:
in this thesis, such term is used to denote a standardized framework that supports the modeling, deployment, management, and interoperability of ecosystems of Digital Twins throughout their lifecycle.
%
Rather than treating data management as a supporting component of the platform, this perspective places data at the core of its design, providing a foundation for addressing the architectural requirements of interoperable digital twin ecosystems.
%
Accordingly, the thesis investigates the Digital Twin Platform not merely as an infrastructure for hosting digital twins, but as a software and data management foundation through which heterogeneous digital twins can be integrated into interoperable ecosystems.

\subsection*{Problem Statement}

The goal of this thesis is to contribute to the development of interoperable Digital Twin ecosystems through the investigation of architectures and methods for Digital Twin Platforms, adopting a data-oriented perspective.
The research focuses on understanding the data management challenges arising from Digital Twin applications and on identifying the models, technologies, methods, and architectural foundations required to address them within a common platform.
%
Working towards this goal, the thesis investigates the following research questions:

\begin{rqlist}
\item[RQ1] \textit{What are the data characteristics and requirements emerging from Digital Twin applications, and to what extent can existing data management technologies satisfy them?}

\item[RQ2] \textit{How can data models, data management approaches, and governance mechanisms support the integration and interoperability of heterogeneous Digital Twin data?} 

\item[RQ3] \textit{How can software architectures support the deployment and management of heterogeneous Digital Twins while satisfying their data and operational requirements?} 

\end{rqlist}


Together, these questions follow the bottom-up perspective adopted throughout the thesis, moving from the characteristics and requirements of Digital Twin data toward the data models, management mechanisms, governance approaches, and software architectures required to support interoperable Digital Twin ecosystems. The answers to these questions contribute to the definition of the architectural and methodological foundations of a Digital Twin Platform, conceived as a common environment for the deployment, management, integration, and interoperability of ecosystems of Digital Twins.

\subsection*{Contributions}
\mpa{list of publications with brief description}

\subsection*{Structure of the thesis}
Accordingly, the thesis is organized into three main parts, each addressing a different aspect of the research problem and contributing to the development of the Digital Twin Platform from complementary perspectives.

\paragraph{Foundations}
\Cref{part:foundations} provides the background notions and concepts required to understand the context and research problem addressed in this thesis.
%
In \Cref{chap:evolving_intertwined_world}, a background is provided on the growing value of data in the digital society, illustrating how the increasing availability of large-scale data has enabled the adoption of advanced data-intensive technologies, including digital twins, and how data management solutions are evolving to meet these demands.
%
In \Cref{chap:digital_twins} the thesis provides a background on digital twins, illustrating their evolution from a research concept to an adopted technological paradigm, and highlighting the challenges that currently hinder their large-scale adoption.
%
Finally, in \Cref{chap:data_oriented_perspective} and within the broader vision of a Digital Twin Platform, the thesis identifies the main characteristics and requirements of digital twin data, illustrating how they can be used to inform the design of data models, management mechanisms, and software infrastructures capable of supporting interoperable digital twin ecosystems.

\paragraph{Architectures}
In \Cref{part:architectures}, the thesis investigates the architectural foundations required to support interoperable digital twin ecosystems within a Digital Twin Platform.
%
Specifically, in \Cref{chap:dataplat} and within a national research project, the thesis explores a data platform solution tailored for precision agriculture, identifying and illustrating the requirements of digital twin applications and how they can be addressed through a general-purpose software infrastructure.
%
This work provides a concrete architectural perspective on how the data and operational requirements identified in the foundations of the thesis can be translated into platform capabilities.
%
Complementarily, following the requirements identified in \Cref{chap:dataplat}, \Cref{chap:stgraph} presents a data model for digital twins, illustrating how such requirements can be addressed through a dedicated and novel data management approach within a Digital Twin Platform.

\paragraph{Methods and applications}
In \Cref{part:methods-applications} and based on the previous findings, the thesis investigates governance techniques and methods facilitating the management of digital twin data within a Digital Twin Platform.
%
In \Cref{sec:process-driven-design}, a methodology is proposed to guide the deployment of data-intensive applications within a data platform, illustrating how the characteristics of the data can be used to produce a blueprint for the design of the underlying software infrastructure.
%
Similarly, in \Cref{sec:llm-guided-metadata-exploration}, a method is proposed to facilitate the governance of heterogeneous data, leveraging large language models to support the discovery and understanding of their metadata.
%
These methods complement the architectural and data management foundations of the Digital Twin Platform by addressing the processes required to design, deploy, and govern data-intensive applications within it.

Finally, in \Cref{chap:applications}, the thesis illustrates some applications relying on the proposed Digital Twin Platform.
%
Within \Cref{sec:smarter}, the thesis presents a digital twin application for precision agriculture and illustrates how the proposed platform can be used to support its deployment and operation.
%
In \Cref{sec:s2g}, the thesis presents a novel technique for anomaly detection in complex sensor networks, illustrating how, within a Digital Twin Platform, the previously identified data model and management techniques can enable novel data analytics.
